\chapter[不曾忘却]{不曾忘却}

\chaptersubtitle{《魔法少女和美》}

《魔圆》电视动画在展开过程中确立了三项基本设定，而这三者结合起来，也给本部分要讨论的三部外传作品出了一道难题。
第一，自魔法少女诞生之初起，魔法少女便是魔女，而且始终如此。
魔女这一形象原本体现的是父权制对女性力量的恐惧\cite{gilbert1979snow-white}。
因此，《魔圆》中的魔法少女一方面天然构成了对支配性权力结构的挑战，另一方面又很容易被解读为一种既会自身腐化，又会带来腐化的因素。
第二，《魔圆》的世界由孵化者的系统支配：
这套系统通过收集魔法少女堕至魔女时的情感衰败，抵消宇宙的熵增，而这一切并未征得魔法少女本人同意。
换言之，这套支配性权力结构以女性为剥削对象，也对女性极端不公。
第三，这套权力结构高度稳定，直到小圆许下愿望、创造出新的系统时才遭到推翻。

姑且搁下小圆的系统是否比丘比的系统更公正这一问题（不过可参见间奏二与第十七章），第三项设定意味着：
凡是以动画主时间线为背景，或以丘比系统占据支配地位的任一条“先前”时间线为背景的外传作品，最终都不得不把丘比的系统当作无从逃避、无从改变的既定事实。
再考虑到丘比的系统对魔法少女力量的限制与剥削（第二项设定），以及这种力量又与一个把女性力量描写成腐化且危险之物的传统联系在一起（第一项设定），外传作品便很容易落入这样的陷阱：
把魔法少女写成危险的存在，再把丘比的系统说成遏制她们所必需的恶，甚至认为这套系统总体上利大于弊。

漫画《魔法少女和美》（平松正树编剧，天杉贵志作画）落入的恰恰就是这个陷阱。

《魔法少女和美》开篇时，主人公和美几乎是一张彻底的白纸：
她赤身裸体，被关在一只行李箱中，对自己的过去毫无记忆。
这是那套颇具“恶名”的\emph{Turkey City Lexicon}\footnote{%
    \emph{Turkey City Lexicon} 是一部针对科幻文学的专门批评术语表，由写作工作坊发展形成，并在各工坊间流传。
}（《土耳其城写作术语表\footnote{%
    译注：
    该术语表意在给小说中那些反复出现，但很难用一句话说清的写作通病起名字，方便作者和评审直接讨论。
}》）\cite{shiner-sterling2009turkey-city}中所谓“白房间综合征”的一种极端形式：
故事开头只有一名角色独处白色房间，困惑于自己为何会身在此处。
这部词典颇为刻薄，但或许也准确地指出，作者这样开篇是因为还没有想清楚故事背景和人物设定。
这番批评是否适用于《魔法少女和美》，尚难断言。
一方面，和美的失忆确实承担着重要的情节功能：
读者得以借她之眼，逐步发现她的真实身份以及昴宿圣团的真相。
不过，《魔法少女和美》早在电视动画完结之前便已开始连载，而且时间并不短：
按月连载，共有二十三话。
漫画后期揭晓的许多真相，在逻辑上都有赖于《魔圆》临近结尾时才揭晓的信息。
我们目前并不清楚，这些真相究竟是为了避免外传漫画提前泄露本篇剧情，才被刻意留到后来，还是因为平松与天杉起初并不知道相关剧情，不得不在连载途中调整方向，以免与本篇相矛盾。

随后，漫画揭示了和美与朋友们的魔法少女身份：
她们隶属一支名为“昴宿圣团”的七人小队，协助她们的则是一个看似更加友善的丘比变体——柔比。
随着情节展开，读者逐渐发现，昴宿圣团远比初见时显得阴暗。
最终我们得知，她们原本是一群并肩合作的魔法少女，和美则是她们的导师。
然而，和美后来变成了魔女。
众人合力让和美恢复为近似人类的形态，又杀死当地的丘比个体，将其改造成柔比。
与此同时，她们还用魔法笼罩整座城市，使接替它的柔比无法被察觉。
昴宿圣团误以为柔比能够替她们净化灵魂宝石，实际得到的却只是一种净化的假象。
在这层假象之下，她们仍在缓慢地走向魔女化。
而此时的和美，也不过是复活和纱美千瑠的一连串尝试中最新的产物。
每一个复活体最终都会崩坏，变成魔女。
到了结局，昴宿圣团除两人外其余全部丧生，一个巨大的聚合魔女也随之诞生。
但和美终于意识到，自己已是不同于和纱美千瑠的另一个人，于是与丘比订立契约，成为魔法少女，拯救了城市与幸存的圣团成员。

这个故事本就相当曲折，背景情节又是依照和美发现真相的顺序，而非事件实际发生的顺序讲述，因而显得更加复杂。
它所做的，是把我在第九章讨论过的魔女与魔法少女的一体性落实为纸面上的事实。
然而，它并没有借此为魔女正名。
相反，它只是把魔法少女也写成了某种魔女。
这种处理在视觉上体现为和美两套装束的鲜明对比：
一套是我们姑且可以称作“日常”装束的衣服，在漫画第二话中称其“更适合年龄小得多的孩子”；
另一套则是她的魔法少女服装，不仅比电视动画中的服装暴露得多，也更像是经过性化处理的万圣节女巫装——
宽檐高顶的黑帽、超短的黑裙、裸露的腰腹和紧身上衣。
反观小圆饰有褶边的粉色连衣裙、沙耶香颇具少年感的铠甲，乃至电视动画中最为性感化的装束——
麻美高领束腰的连衣裙与打底裤，或杏子拉链几乎全敞开的运动衫与穿在长裤外的短裙——
漫画在天真孩童与强大、危险而带有性意味的魔法少女之间，划出了清晰可见的界线。

这正是漫画与电视动画最为鲜明的差异。
在《魔圆》中，魔法少女与魔女的一体性服务于作品对希望与绝望循环的论述。
但到了《魔法少女和美》中，它却变成了对圣母／荡妇情结的接受。
本作中的魔法少女无不藏有可怕的秘密：
要么她们本身就是魔女（和美），要么她们过去曾做过骇人的事（昴宿圣团）。
昴宿圣团的成员无一例外，都受困于成为魔法少女之前犯下的过错。
她们都曾欺骗和美，也都曾协助猎捕、囚禁其他魔法少女，以免她们变成魔女。
这里的魔女化不再像《魔圆》中沙耶香的经历那样，是一步步堕入绝望的结果，而只是魔法使用过度之后便会发生的事。
换言之，魔女不再诞生于绝望，而成了女性行使力量的结果。
更露骨的是，故事前期的反派本身就是魔法少女，并且能够直接把女性变成伪魔女。

当真正的魔女终于现身时，同样与《魔圆》中的形象大相径庭。
《魔圆》暗示魔女没有理智、受困于妄想，也无法离开自己的结界，但《魔法少女和美》中的魔女却能直接接触并感知世界，其中一部分甚至具有人格与智力。
因此，她们不再是对《浮士德》的指涉，也不再是女性绝望的血肉化身，而只是魔法少女躯体的怪诞变异。
归根结底，在《魔法少女和美》中，魔法少女与魔女的唯一区别，不过是前者“可爱”或“性感”，后者则身体畸变或呈怪物之形。
换言之，魔女的丑陋并不在于她与负面情感的联系，而在于她已不再具有供男性凝视的吸引力。
而这部漫画诉诸男性凝视的频率，远远高于电视动画。

结局同样否定了《魔圆》所呈现的魔女本质。
当和美与两名幸存的朋友重新回到丘比的系统之中后，漫画借和美之口直截了当地断言：
她们不会变成魔女。
《魔圆》中的希望与绝望的循环由此遭到彻底否定。
取而代之的，是这样一种对女性力量的描写：
女性力量只有受到约束、仅用于对抗自身的负面表现，才被允许继续存在。
漫画宣称，只要魔法少女把自己局限于同魔女战斗之中，不去挑战丘比及其腐化的系统，她们便可以继续生存、继续行使自己的力量。

归根结底，《魔法少女和美》精准地落入了这个陷阱：
它声称，必须依靠丘比的系统来约束女性，防止她们以可能挑战现状的方式行使力量，因此这套系统不可或缺。但愿这只是无心之失。
